% =====================================================================
\chapter{F2: the CRT group key leaks -- to members and to outsiders}\label{ch:f2}
% =====================================================================
\begin{flawbox}[title={Flaw F2 -- the broadcast polynomial reveals the decrypting keys}]
The polynomial $\Lambda(x)=\kappa_{tu}\sum_i\omega_i\varphi_i+\prod_i(x-\sigma_i)$ hides only its \emph{constant} term. Every other coefficient is a
symmetric function of the members' decrypting keys, published in the clear.
\begin{itemize}[leftmargin=*]
\item \textbf{F2(a) -- revocation fails.} Any member recovers \emph{every} $\sigma_i$; after being removed it still computes each new group key.
\item \textbf{F2(b) -- backward secrecy fails.} A UAV that joins recovers the old members' $\sigma_i$ and opens recorded old group keys.
\item \textbf{F2(c) -- a passive outsider wins.} An eavesdropper who never registered records the two broadcasts around one join or leave and recovers
every $\sigma_i$ and \emph{both} group keys.
\end{itemize}
\end{flawbox}

\section{One value is the same for everybody}
\begin{lemma}[Common value]\label{lem:common}
Let $C=\kappa_{tu}\sum_{i=1}^n\omega_i\varphi_i$. Then $\Lambda(\sigma_j)=C$ for every member $j$.
\end{lemma}
\begin{proof}
At $x=\sigma_j$ the product $\prod_i(x-\sigma_i)$ contains the factor $(\sigma_j-\sigma_j)=0$, so only the constant $C$ remains.
\end{proof}
\noindent The paper uses this fact to deliver the key ($C\bmod\sigma_j=\kappa_{tu}$, Lemma~\ref{lem:gd6}). But the same fact means that every member
knows the \emph{mask} $C$ itself, not only $C\bmod\sigma_j$.

\section{The insider attack: F2(a) and F2(b)}
\begin{theorem}[Insider key recovery]\label{thm:insider}
A member $j$ that knows only its own $\sigma_j$ and the broadcast coefficients $\{a_k\}$ recovers every member's decrypting key in polynomial time.
\end{theorem}
\begin{proof}
By Lemma~\ref{lem:common} member $j$ computes $C=\Lambda(\sigma_j)$, so it knows
\[
\boxed{\;\Lambda(x)-C=\prod_{i=1}^{n}(x-\sigma_i)\;}
\]
-- the broadcast polynomial with $a_0$ replaced by $a_0-C$. This is a monic integer polynomial whose roots are exactly the $\sigma_i$; all of them are
integers (Remark~\ref{rem:typing}), so by the Rational Root Theorem (\ref{thm:rrt}) and Proposition~\ref{prop:roots} they are found in time polynomial in
$n$ and the bit length.
\end{proof}

\begin{corollary}[F2(a): revocation fails]\label{cor:revoke}
Let member $j$ be removed by DU1 and let $k$ be any member that stays. From the broadcast before its removal, $j$ knows $\sigma_k$
(Theorem~\ref{thm:insider}); DU2 keeps $\sigma_k$ unchanged; hence $j$ computes the new key $\kappa^*_{tu}=\Lambda^*(\sigma_k)\bmod\sigma_k$.
\end{corollary}

\begin{corollary}[F2(b): backward secrecy fails]\label{cor:backward}
Let UAV $k$ join in DU1. Applying Theorem~\ref{thm:insider} to $\Lambda^*$, it learns the decrypting keys of the old members, which DU2 kept unchanged. For
any recorded earlier broadcast $\Lambda$ and any old member $j$, it computes the old key $\kappa_{tu}=\Lambda(\sigma_j)\bmod\sigma_j$.
\end{corollary}

\begin{figure}[H]
\centering
\resizebox{\linewidth}{!}{%
\begin{tikzpicture}[font=\sffamily\small,node distance=0.45cm]
  \node[fio=bad,text width=2.5cm] (in) {receive $\{a_k\}$ (GD4) as a member $j$};
  \node[fproc=bad,right=of in,text width=2.4cm] (c) {$C=\Lambda(\sigma_j)$\\ (Lemma~\ref{lem:common})};
  \node[fproc=bad,right=of c,text width=2.7cm] (p) {$\Lambda(x)-C=\prod_i(x-\sigma_i)$};
  \node[fproc=vehc,right=of p,text width=2.5cm] (r) {integer roots\\ (Prop.~\ref{prop:roots})\\ $\Rightarrow$ all $\sigma_i$};
  \node[fdec=bad,right=of r] (d) {later\\ event?};
  \node[fstart=bad,above right=0.2cm and 0.7cm of d,text width=3.4cm] (a) {\textbf{F2(a)} removed: $\kappa^*=\Lambda^*(\sigma_k)\bmod\sigma_k$};
  \node[fstart=bad,below right=0.2cm and 0.7cm of d,text width=3.4cm] (b) {\textbf{F2(b)} joined: old $\kappa=\Lambda(\sigma_j)\bmod\sigma_j$};
  \draw[flow] (in) -- (c); \draw[flow] (c) -- (p); \draw[flow] (p) -- (r); \draw[flow] (r) -- (d);
  \draw[flow] (d) |- node[lbl,pos=0.3,left]{revoked} (a); \draw[flow] (d) |- node[lbl,pos=0.3,left]{newcomer} (b);
\end{tikzpicture}}
\caption{Flowchart of the insider attack (Theorem~\ref{thm:insider}). A member evaluates the broadcast polynomial at its own key, subtracts the result and finds the roots; this gives every member's key. Depending on what happens later, the same member can break revocation, F2(a), or backward secrecy, F2(b).}
\label{fig:insiderflow}
\end{figure}



\section{The passive outsider attack: F2(c)}
The insider attack needs one $\sigma_j$. The outsider attack needs none: it uses only the public coefficients of two broadcasts.

\begin{lemma}[Vieta sum leak]\label{lem:sum}
If $n\ge2$, the coefficient $a_{n-1}$ of $\Lambda(x)$ equals $-(\sigma_1+\dots+\sigma_n)$. More generally $a_k$ for $1\le k\le n-1$ are the coefficients of
$\prod_i(x-\sigma_i)$ exactly, without any mask.
\end{lemma}
\begin{proof}
The mask $C$ is a constant, so it changes only $a_0$. The remaining coefficients are those of $\prod_i(x-\sigma_i)$, given by Vieta's formulas
(Theorem~\ref{thm:vieta}); in particular $a_{n-1}=-e_1(\sigma)=-\sum_i\sigma_i$.
\end{proof}

\begin{theorem}[Passive outsider key recovery]\label{thm:outsider}
Let $\Lambda$ and $\Lambda'$ be two consecutive broadcasts whose member sets differ by exactly one UAV with key $\sigma$, and suppose at least two members
are common to both sets. An eavesdropper that knows only $\{a_k\}$ and $\{a'_k\}$ computes, in polynomial time, $\sigma$, every other decrypting key, the
old group key $\kappa_{tu}$ and the new group key $\kappa^*_{tu}$.
\end{theorem}
\begin{proof}
Write $p(x)=\prod_{\text{old}}(x-\sigma_i)=\sum p_kx^k$ and $q(x)=\prod_{\text{new}}(x-\sigma_i)=\sum q_kx^k$. By Lemma~\ref{lem:sum} the eavesdropper
knows every coefficient of $p$ and $q$ except the constants $p_0$ and $q_0$.

\emph{Leave} ($\deg q=\deg p-1$). Then $p(x)=(x-\sigma)\,q(x)$. The sums are public, so $\sigma=-p_{n-1}+q_{n-2}$. Comparing the coefficients of $x^1$
in $p=(x-\sigma)q$ gives $p_1=q_0-\sigma q_1$, hence
\[
\boxed{\;q_0=p_1+\sigma\,q_1\;}
\]
where $p_1=a_1$ and $q_1=a'_1$ are public (because at least two members remain, $q_1$ is not the masked constant). Now $q$ is known completely, its roots
are the remaining keys (Proposition~\ref{prop:roots}), and $p=(x-\sigma)q$ is known too. With any key, the attacker evaluates $\Lambda(\sigma_j)\bmod\sigma_j$
and $\Lambda'(\sigma_j)\bmod\sigma_j$ like a member and obtains both group keys.

\emph{Join} ($\deg q=\deg p+1$). Now $q(x)=(x-\sigma)\,p(x)$, $\sigma=-q_{n}+p_{n-1}$, and the $x^1$ coefficients give
$\boxed{\,p_0=q_1+\sigma\,p_1\,}$; the rest is identical with the roles of $p$ and $q$ swapped.
\end{proof}

\begin{algorithm}{OutsiderRecover$(\{a_k\},\{a'_k\})$ -- two recorded broadcasts}\label{alg:outsider}
\begin{enumerate}[leftmargin=*,label=\arabic*.]
\item $n\leftarrow\deg\Lambda$, \;$n'\leftarrow\deg\Lambda'$. If $|n-n'|\neq1$, stop (outside the theorem).
\item Sums: $S=-a_{n-1}$, \;$S'=-a'_{n'-1}$; \;$\sigma\leftarrow|S-S'|$.
\item If $n'=n-1$ (leave): $q_0\leftarrow a_1+\sigma\,a'_1$; \;$f(x)\leftarrow x^{n'}+a'_{n'-1}x^{n'-1}+\dots+a'_1x+q_0$ (the new product).\\
If $n'=n+1$ (join): $p_0\leftarrow a'_1+\sigma\,a_1$; \;$f(x)\leftarrow x^{n}+a_{n-1}x^{n-1}+\dots+a_1x+p_0$ (the old product).
\item Find the integer roots of $f$ (Proposition~\ref{prop:roots}); together with $\sigma$ these are all keys of both sets.
\item Output $\kappa_{tu}=\Lambda(\sigma_j)\bmod\sigma_j$ and $\kappa^*_{tu}=\Lambda'(\sigma_k)\bmod\sigma_k$ for any old $j$ and new $k$.
\end{enumerate}
\end{algorithm}

\begin{figure}[H]
\centering
\resizebox{\linewidth}{!}{%
\begin{tikzpicture}[font=\sffamily\small]
  \node[fio=deep,text width=2.4cm] (b0) at (0,1.1) {record broadcast\\ $\{a_k\}$ (before)};
  \node[fio=deep,text width=2.4cm] (b1) at (0,-1.1) {record broadcast\\ $\{a'_k\}$ (after)};
  \node[fproc=bad,text width=2.9cm] (s) at (4.4,0) {sums $S=-a_{n-1}$,\\ $S'=-a'_{n'-1}$;\\ $\sigma=|S-S'|$};
  \node[fdec=bad] (d) at (8.0,0) {leave or\\ join?};
  \node[fproc=bad,text width=2.6cm] (l) at (11.6,1.3) {leave: $q_0=a_1+\sigma a'_1$};
  \node[fproc=bad,text width=2.6cm] (j) at (11.6,-1.3) {join: $p_0=a'_1+\sigma a_1$};
  \node[fproc=vehc,text width=2.6cm] (r) at (15.2,0) {full product\\ $\Rightarrow$ integer roots\\ $\Rightarrow$ all $\sigma_i$};
  \node[fstart=bad,text width=2.4cm] (o) at (18.6,0) {\textbf{old and new}\\ group keys};
  \draw[flow] (b0.east) -- (s.west); \draw[flow] (b1.east) -- (s.west); \draw[flow] (s) -- (d);
  \draw[flow] (d.north) |- node[lbl,pos=0.3,left]{$n'=n-1$} (l.west); \draw[flow] (d.south) |- node[lbl,pos=0.3,left]{$n'=n+1$} (j.west);
  \draw[flow] (l.east) -| (r.north); \draw[flow] (j.east) -| (r.south); \draw[flow] (r) -- (o);
\end{tikzpicture}}
\caption{Flowchart of the passive outsider attack (Algorithm~\ref{alg:outsider}). The attacker records the coefficient sets before and after one membership change, obtains the changed member's key from the two coefficient sums, removes the mask, and finds all keys and both group keys. It needs no registration and no secret.}
\label{fig:outflow}
\end{figure}


\subsection*{Several UAVs leaving (or joining) at once}
DU1 is a \emph{batch} update: $\ominus$ UAVs may leave or $\oplus$ may join in one step. The same algebra extends.

\begin{proposition}[Batch changes]\label{prop:batch}
Suppose $m\ge1$ UAVs leave in one update and $r\ge m+1$ members remain (or $m$ UAVs join a group of $r\ge m+1$). Then the two broadcasts determine all
keys of both groups in polynomial time.
\end{proposition}
\begin{proof}
For a leave write $p=q\cdot d$ with $q=\prod_{\text{stay}}(x-\sigma_i)$ of degree $r$ and $d=\prod_{\text{leave}}(x-\sigma_i)=x^m+d_{m-1}x^{m-1}+\dots+d_0$.
Comparing the coefficients of $x^{n-j}$ for $j=1,\dots,m$ (with $n=r+m$) gives
\[
d_{m-j}=p_{n-j}-\sum_{i=1}^{j}q_{r-i}\,d_{m-j+i},
\]
which uses only $p_{n-j}$ ($n-j\ge r\ge1$, unmasked), $q_{r-i}$ with $r-i\ge r-m\ge1$ (unmasked) and already computed $d$'s. So $d$ is found top-down.
The coefficient of $x^m$ then gives $q_0=p_m-\sum_{i=1}^{\min(m,r)}q_i\,d_{m-i}$. Now $d$ and $q$ are known; their roots are all keys. A batch join is the
same with the roles of the two broadcasts swapped.
\end{proof}


\begin{remark}[Scope]
Theorem~\ref{thm:outsider} is proved for one join or one leave between two broadcasts, and Proposition~\ref{prop:batch} for batches of pure leaves or pure
joins with more members staying than changing. Two cases remain open: a \emph{single} broadcast with no
second one to compare, and one update that mixes joins and leaves. We also confirmed Theorem~\ref{thm:outsider} numerically at the paper's own size
(160-bit $\sigma_i$, 12 UAVs, one leave and one join): every $\sigma_i$ and both group keys were recovered.
\end{remark}

\section{How this contradicts the paper}
\begin{claimbox}
``the dynamic key batch updating process can be securely executed in the open wireless environment since the malicious entities that successfully
collect the coefficient set $\{a_i\}_{i\in[0,n]}$ can not acquire the $\kappa_{tu}$ without the correct decrypting key $\sigma_i$.'' \hfill (Section IV-D)
\end{claimbox}
Theorem~\ref{thm:outsider} refutes this sentence directly: an entity that collects two coefficient sets, and nothing else, acquires $\kappa_{tu}$
\emph{and} every $\sigma_i$. Corollaries~\ref{cor:revoke} and~\ref{cor:backward} refute ``supports multi-device revocation and new device participation''
(Section~IV-D) and the proof of Theorem~5 (``if the $\sigma_i$ has not been removed by TUAV'').

\section{Root cause, and why a small patch is not enough}
By Vieta's formulas the scheme \emph{publishes} the elementary symmetric functions $e_1,\dots,e_{n-1}$ of the members' secrets and masks only
$e_n$. A set whose symmetric functions are public is not hidden: one member pins it down completely (Theorem~\ref{thm:insider}), and two broadcasts that
differ by one member reveal it to anyone (Theorem~\ref{thm:outsider}).

Two patches come to mind, and neither is attractive. (i) Giving every member a \emph{fresh} $\sigma_i$ at each update would stop the leaks across updates,
but $\sigma_i=h_1(s_i\xi_i)P$ changes only with a new $\xi_i$, i.e.\ a new US/GA round through the TA for \emph{every} member at \emph{every} update --
exactly the cost DU was meant to avoid. (ii) Masking all coefficients would require each member to remove the mask, i.e.\ a secret shared with each
member -- which is per-member key delivery. We therefore drop the public polynomial entirely (C1).

\begin{remark}[Size of the broadcast]\label{rem:size}
The coefficient $a_k$ is a symmetric function of degree $n-k$ in 160-bit values, so it has about $160(n-k)$ bits, and one coefficient set has about
$160\big(\tfrac{n(n-1)}{2}+n+1\big)$ bits $\approx10n^2$ bytes. GD3 sends it ``one by one to each vehicle'', so one key distribution moves about
$10n^3$ bytes: roughly $17.4$\,MB for $n=120$.
\end{remark}


\begin{limitbox}
All results use the integer reading of $\sigma_i$ (Remark~\ref{rem:typing}), which is the only reading under which GD6 is defined. CRT-based group-key
schemes are already known to be fragile when a user is \emph{corrupted}~\cite{sp24}; our insider attack is an explicit instance of that weakness for this
scheme. The passive \emph{outsider} attack needs no corruption at all, and we have not found it reported for this scheme or its 94 citing works.
\end{limitbox}

\begin{fixbox}[title={How C1 removes F2 (Chapter~\ref*{ch:c1})}]
No function of the members' secrets is ever published. The group key travels as one AEAD ciphertext per member, each under its own fresh
Diffie--Hellman key, and every join or leave triggers a fresh, independent key (Theorem~\ref{thm:gksecrecy}).
\end{fixbox}
